%==============================================================================
% CAPÍTULO 12 — SDD: DESENVOLVIMENTO ORIENTADO POR ESPECIFICAÇÃO
% PARTE III — Engenharia de Software Odontológico com SDD e TDD
%==============================================================================
\chapter{SDD: Desenvolvimento Orientado por Especificação}
\label{cap:sdd-metodologia}
\niveltres

Este capítulo inaugura a Parte~III do livro, dedicada à engenharia de
software odontológico com metodologias formais de especificação e teste.
Aqui, estabelecemos o alicerce conceitual e prático do
\textit{Specification-Driven Development} (SDD) — a disciplina de nunca
implementar sem antes especificar. Em um domínio onde erros de software
podem significar diagnósticos equivocados com consequências clínicas
graves, a especificação formal não é luxo acadêmico: é exigência
ética e profissional.

\section{O que é Specification-Driven Development?}

O \textit{Specification-Driven Development} (SDD) é uma metodologia de
desenvolvimento de software que estabelece que \textbf{toda implementação
deve ser precedida por uma especificação formal verificável}. Diferentemente
do desenvolvimento \textit{ad hoc} — em que o programador ``sai codando''
— o SDD impõe um contrato explícito entre a intenção do desenvolvedor e
o comportamento do sistema.

\indice{Specification-Driven Development}

A origem intelectual do SDD remonta à engenharia de software formal da
década de 1980, quando \citeonline{beck2003test} já argumentavam que
especificações escritas antes do código reduzem defeitos em até 60\%.
No entanto, o SDD moderno distingue-se por três características:

\begin{enumerate}
  \item \textbf{Especificações executáveis}: não são documentos estáticos,
    mas artefatos que podem ser validados automaticamente contra a
    implementação.
  \item \textbf{Critérios de aceitação mensuráveis}: cada especificação
    contém thresholds numéricos verificáveis (ex.: ``acurácia $>$ 0,88'').
  \item \textbf{Rastreabilidade bidirecional}: é possível navegar da
    especificação ao código e do código à especificação.
\end{enumerate}

\notaexplicativa{O termo \textit{specification-driven} não deve ser confundido
com \textit{specification-based testing}, que é uma técnica de caixa-preta.
O SDD é uma metodologia completa de desenvolvimento, enquanto
specification-based testing é apenas uma das técnicas de teste que podem
ser usadas dentro do SDD.}

\subsection{SDD, TDD e BDD: um triângulo complementar}

É comum que o leitor, ao ouvir falar de ``especificação antes do código'',
pense imediatamente em TDD (\textit{Test-Driven Development}) ou BDD
(\textit{Behavior-Driven Development}). A Tabela~\ref{tab:sdd-tdd-bdd}
esclarece as diferenças e complementaridades.

\begin{table}[H]
  \centering
  \caption{Comparação entre SDD, TDD e BDD.}
  \label{tab:sdd-tdd-bdd}
  \begin{tabularx}{\textwidth}{lXXX}
    \toprule
    \textbf{Aspecto} & \textbf{SDD} & \textbf{TDD} & \textbf{BDD} \\
    \midrule
    Foco & Especificação formal & Teste automatizado & Comportamento \\
    Artefato principal & Documento de spec & Caso de teste & Cenário Gherkin \\
    Quem escreve & Arquiteto / Especialista & Desenvolvedor & QA / Product Owner \\
    Verificabilidade & Automática (SpecVerifier) & Automática (pytest) & Automática (Cucumber) \\
    Nível de abstração & Sistema/função & Unidade/integração & Funcionalidade \\
    Exemplo odonto & ``accuracy $>$ 0.88'' & ``assert acc $>$ 0.88'' & ``Dado uma radiografia...'' \\
    \bottomrule
  \end{tabularx}
\end{table}

Na prática, as três metodologias se complementam: o SDD define \textbf{o quê}
deve ser construído; o TDD garante \textbf{como} verificar; o BDD traduz
\textbf{para quem}. Neste livro, adotamos SDD como metodologia-mãe, com
TDD como prática de verificação contínua.


% === FIGURAS ADICIONADAS (OdontoIA Dark) ===
% ======================================================================
% FIGURA: Pipeline SDD+TDD Odontológico (Capítulo 12)
% ======================================================================
\begin{figure}[H]
\centering
\begin{tikzpicture}[
    node distance=1.2cm and 2.5cm,
    box/.style={rectangle, draw=blueaccent, fill=blueaccent!15, align=center, minimum width=2.8cm, align=center, minimum height=0.9cm, font=\footnotesize, rounded corners=3pt, align=center},
    testbox/.style={rectangle, draw=codegreen, fill=codegreen!15, align=center, minimum width=2.8cm, align=center, minimum height=0.9cm, font=\footnotesize, rounded corners=3pt, align=center},
    arrow/.style={-latex, thick},
]
% Fase 1: SDD
\node[align=center, box, fill=blueaccent!25] (spec) {\textbf{ESPECIFICAR}\\Problema clínico\\Entrada/Saída\\Critérios aceitação};
\node[align=center, box, right=of spec] (data) {Preparar\\Dados\\Anotar\\Validar};

% Fase 2: TDD
\node[align=center, testbox, below=of spec] (red) {\textbf{RED}\\Escrever teste\\Ver falhar};
\node[align=center, box, right=of red] (green) {\textbf{GREEN}\\Implementar\\mínimo código};
\node[align=center, testbox, right=of green] (refactor) {\textbf{REFACTOR}\\Otimizar\\Documentar};

% Fase 3: Verificar
\node[align=center, box, fill=orangeaccent!20, below=of red] (validate) {Validar\\Métricas clínicas\\Sens/Espec $\ge 0.85$};
\node[align=center, box, fill=orangeaccent!20, right=of validate] (deploy) {Publicar\\Modelo + Testes\\DOI + Dataset};

% Setas
\draw[arrow] (spec) -- (data);
\draw[arrow] (spec) -- (red);
\draw[arrow] (data) -- (green);
\draw[arrow] (red) -- (green) node[midway, above, font=\tiny] {falha};
\draw[arrow] (green) -- (refactor) node[midway, above, font=\tiny] {passa};
\draw[arrow] (refactor) -- (validate);
\draw[arrow] (validate) -- (deploy);

% Ciclo
\draw[arrow, bend left=40, dashed, gray] (refactor.north) to (red.north) node[midway, above, font=\tiny\color{codegray}] {ciclo};

\node[font=\small\bfseries, above=0.5cm of spec] {SDD + TDD Odontológico};

\end{tikzpicture}
\caption{Pipeline de desenvolvimento orientado por especificação (SDD) e testes (TDD) para modelos de IA odontológica. O ciclo RED$\rightarrow$GREEN$\rightarrow$REFACTOR garante que cada função clínica seja testada antes da implementação, reduzindo erros em diagnóstico automatizado.}
\label{fig:sdd-tdd-pipeline}
\end{figure}


\section{Por que SDD na Odontologia?}

A Odontologia é um domínio de \textbf{segurança crítica} (\textit{safety-critical})
no sentido de que erros de software podem causar danos a pacientes. Um falso
negativo em um detector de cáries pode atrasar o tratamento; um falso positivo
em um classificador de lesões pode gerar biópsias desnecessárias.

\indice{segurança clínica} \indice{software safety-critical}

\citacaoativa{doi-1016-j-jdent-2021-103830}{10.1016/j.jdent.2021.103830}{%
  ``A implementação clínica de sistemas de IA em Odontologia requer
  protocolos rigorosos de validação, especificação formal de requisitos
  clínicos e testes de desempenho em cenários realísticos antes de
  qualquer uso assistencial.''}

Este trecho, extraído de Schwendicke et al.\ (2021), sintetiza a
motivação central para o SDD odontológico. Vejamos os argumentos
específicos:

\subsection{Segurança clínica como requisito não negociável}

Na aviação civil, o padrão DO-178C exige rastreabilidade completa entre
requisitos e código para sistemas críticos de voo. Na Odontologia, ainda
não temos um equivalente regulatório — o que torna ainda \textbf{mais}
importante que o desenvolvedor adote voluntariamente práticas rigorosas.

A Tabela~\ref{tab:safety-comparison} compara domínios \textit{safety-critical}.

\begin{table}[H]
  \centering
  \caption{Comparação de níveis de criticidade entre domínios.}
  \label{tab:safety-comparison}
  \begin{tabularx}{\textwidth}{lccX}
    \toprule
    \textbf{Domínio} & \textbf{Criticidade} & \textbf{Padrão} & \textbf{Consequência de falha} \\
    \midrule
    Aviação & Nível A (catastrófico) & DO-178C & Perda de vidas \\
    Medicina (geral) & Nível B-C (grave) & ISO 13485 / IEC 62304 & Lesão ao paciente \\
    Odontologia IA & Nível B-C (grave) & Não regulado (ainda) & Diagnóstico incorreto \\
    Software empresarial & Nível D-E (menor) & ISO 9001 & Perda financeira \\
    \bottomrule
  \end{tabularx}
\end{table}

\destaque{O software odontológico com IA opera em um vácuo regulatório.
Enquanto a ANVISA e o FDA discutem frameworks para \textit{SaMD}
(Software as Medical Device), o desenvolvedor responsável deve adotar
SDD e TDD como autocertificação de qualidade.}

\subsection{Reprodutibilidade como exigência científica}

A crise de reprodutibilidade na ciência — amplamente documentada por
\citeonline{baker2016reproducibility} na \textit{Nature} — afeta também
a IA odontológica. Sem especificações formais, dois grupos de pesquisa
podem implementar ``o mesmo'' classificador de cáries e obter resultados
completamente diferentes, simplesmente porque interpretaram diferentemente
o que significa ``pré-processar uma radiografia''.

\citacaoativa{baker2016reproducibility}{10.1038/533452a}{%
  ``Mais de 70\% dos pesquisadores já tentaram e falharam em reproduzir
  experimentos de outros cientistas, e mais da metade falhou em reproduzir
  seus próprios experimentos.''}

O SDD resolve isso ao tornar \textbf{explícito} o que estava implícito:
o contrato de entrada/saída de cada função, os critérios de aceitação
numéricos e as invariantes do sistema.

\subsection{Rastreabilidade para auditoria e regulamentação}

Quando (não \textit{se}, mas \textit{quando}) a ANVISA e o FDA
estabelecerem requisitos para \textit{SaMD} odontológico, sistemas
desenvolvidos com SDD terão vantagem competitiva: cada decisão de
implementação estará vinculada a uma especificação rastreável, cada
especificação a um requisito clínico, cada requisito a uma evidência
científica.

\section{Escrevendo Especificações: Problema, Entrada, Saída, Critérios}

Uma especificação SDD bem formada possui \textbf{quatro componentes
obrigatórios}:

\begin{enumerate}
  \item \textbf{Problema} (\texttt{problem}): descrição clara do problema
    odontológico que a função resolve.
  \item \textbf{Entrada} (\texttt{input}): tipo, forma e domínio dos
    dados de entrada.
  \item \textbf{Saída} (\texttt{output}): tipo, forma e semântica dos
    dados de saída.
  \item \textbf{Critérios de aceitação} (\texttt{criteria}): lista de
    condições verificáveis, preferencialmente com thresholds numéricos.
\end{enumerate}

\subsection{Anatomia de uma boa especificação}

Considere o seguinte exemplo \textbf{ruim} de especificação:

\begin{mdframed}[linecolor=corTesteFalha, linewidth=2pt, backgroundcolor=corTesteFalha!5]
  \textbf{Especificação VAGA (ruim):} \\
  \textit{``A função deve detectar cáries em imagens com boa acurácia.''}
\end{mdframed}

Esta especificação falha em todos os quatro componentes: o problema é
genérico, a entrada não está tipada, a saída é ambígua e não há critério
verificável.

Agora, a versão \textbf{correta}:

\begin{mdframed}[linecolor=corTesteSucesso, linewidth=2pt, backgroundcolor=corTesteSucesso!5]
  \textbf{Especificação PRECISA (boa):} \\
  \textbf{Problema:} Detectar cáries proximais (ICDAS $\geq$ 2) em radiografias
  bitewing digitais de dentes posteriores permanentes. \\
  \textbf{Entrada:} \texttt{ndarray (224, 224) uint8} — radiografia bitewing
  recortada e normalizada. \\
  \textbf{Saída:} \texttt{dict\{"has\_caries": bool, "confidence": float, "icdas\_score": int\}} \\
  \textbf{Critérios:} accuracy $>$ 0,88; sensitivity $>$ 0,85; specificity $>$ 0,85;
  inference\_time $<$ 200ms.
\end{mdframed}

\notaexplicativa{O tipo \texttt{ndarray} refere-se a um \texttt{numpy.ndarray},
o formato padrão de representação de imagens no ecossistema Python científico.
O sufixo \texttt{uint8} indica inteiros sem sinal de 8 bits (valores 0--255).}

\subsection{Critérios de aceitação clinicamente significativos}

Os critérios de aceitação em SDD odontológico não são métricas arbitrárias:
cada uma tem significado clínico direto. A Tabela~\ref{tab:criterios-clinicos}
mapeia métricas de machine learning para implicações clínicas.

\begin{table}[H]
  \centering
  \caption{Mapeamento de critérios de aceitação para significado clínico.}
  \label{tab:criterios-clinicos}
  \begin{tabularx}{\textwidth}{lXl}
    \toprule
    \textbf{Métrica} & \textbf{Significado clínico} & \textbf{Threshold típico} \\
    \midrule
    Accuracy & Proporção geral de acertos diagnósticos & $>$ 0,88 \\
    Sensitivity & Capacidade de detectar doença quando presente & $>$ 0,85 \\
    Specificity & Capacidade de confirmar ausência de doença & $>$ 0,85 \\
    PPV (Precision) & Probabilidade de doença dado teste positivo & $>$ 0,80 \\
    NPV & Probabilidade de saúde dado teste negativo & $>$ 0,90 \\
    AUC-ROC & Capacidade discriminativa global & $>$ 0,90 \\
    Dice (seg.) & Sobreposição entre segmentação predita e real & $>$ 0,80 \\
    \bottomrule
  \end{tabularx}
\end{table}

\teste{Todos os critérios da Tabela~\ref{tab:criterios-clinicos} são
verificáveis automaticamente via \texttt{odontoia.metrics} e integráveis
ao pipeline de CI/CD.}

\section{O Ciclo SDD: Spec $\rightarrow$ Implement $\rightarrow$ Verify $\rightarrow$ Deploy}

O ciclo de desenvolvimento SDD é iterativo e composto por quatro fases
distintas, como ilustrado na Figura~\ref{fig:sdd-cycle}.

\begin{figure}[H]
  \centering
  \begin{tikzpicture}[
    node distance=2.5cm,
    box/.style={rectangle, draw=corPrimaria, fill=corPrimaria!5, text width=3.5cm, align=center, align=center, minimum height=2cm, rounded corners=4pt},
    arrow/.style={->, >=stealth, thick, corSecundaria},
  ]
    \node[align=center, box] (spec) {\textbf{1. SPEC}\\\footnotesize Escrever especificação\\\footnotesize formal com critérios};
    \node[align=center, box, right=of spec] (impl) {\textbf{2. IMPLEMENT}\\\footnotesize Implementar função\\\footnotesize seguindo a spec};
    \node[align=center, box, below=of impl] (verify) {\textbf{3. VERIFY}\\\footnotesize SpecVerifier + testes\\\footnotesize automatizados};
    \node[align=center, box, below=of spec] (deploy) {\textbf{4. DEPLOY}\\\footnotesize Integrar, documentar\\\footnotesize e publicar};

    \draw[arrow] (spec) -- (impl);
    \draw[arrow] (impl) -- (verify);
    \draw[arrow] (verify) -- (deploy);
    \draw[arrow, dashed, corDestaque] (deploy) -- (spec) node[midway, left, corDestaque] {\footnotesize refinamento};
  \end{tikzpicture}
  \caption{Ciclo de desenvolvimento SDD em quatro fases.}
  \label{fig:sdd-cycle}
\end{figure}

\subsection{Fase 1 — SPEC: Especificação}

A fase de especificação é a mais importante do ciclo: \textbf{se a
especificação estiver errada, todo o resto estará errado}. Nesta fase:

\begin{itemize}
  \item O especialista de domínio (cirurgião-dentista) define o problema
    clínico e os critérios de aceitação clinicamente significativos.
  \item O engenheiro de software traduz esses requisitos para tipos,
    contratos e thresholds numéricos.
  \item A especificação é registrada via decorador \texttt{@spec}.
\end{itemize}

\subsection{Fase 2 — IMPLEMENT: Implementação}

Com a especificação aprovada, o desenvolvedor implementa a função.
\textbf{Nenhuma linha de código é escrita antes que a especificação
esteja aprovada}. A implementação deve:

\begin{itemize}
  \item Respeitar o contrato de tipos definido na especificação.
  \item Incluir docstring explicativa (verificada pelo SpecVerifier).
  \item Ser a implementação \textbf{mínima} que satisfaz os critérios
    (princípio YAGNI: \textit{You Ain't Gonna Need It}).
\end{itemize}

\subsection{Fase 3 — VERIFY: Verificação}

A verificação é automatizada pelo \texttt{SpecVerifier} e pelos testes
TDD (Capítulo~13). O \texttt{SpecVerifier} verifica:

\begin{enumerate}
  \item Se a função possui especificação registrada.
  \item Se os tipos de entrada e saída são compatíveis.
  \item Se a docstring está presente.
  \item Se os critérios são sintaticamente válidos.
\end{enumerate}

\subsection{Fase 4 — DEPLOY: Publicação}

Após verificação bem-sucedida, o código é integrado ao repositório,
documentado e publicado. A especificação viaja junto com o código como
metadado, permitindo auditoria futura.

\section{Templates de Especificação para Problemas Odontológicos}

Para acelerar o processo de especificação, o pacote \texttt{odontoia}
oferece templates pré-configurados para as categorias mais comuns de
problemas odontológicos.

\subsection{Template para Classificação}

\spec{{
\textbf{Template: Classificação Odontológica} \\
\textbf{Problema:} Classificar [tipo de imagem] em [N] categorias clínicas. \\
\textbf{Entrada:} \texttt{ndarray (H, W[, C])} — [descrição da imagem]. \\
\textbf{Saída:} \texttt{dict\{"class": str/int, "probabilities": ndarray, "confidence": float\}} \\
\textbf{Critérios:}
\begin{itemize}
  \item accuracy $>$ [0.85--0.95]
  \item sensitivity $>$ [0.80--0.92]
  \item specificity $>$ [0.80--0.92]
\end{itemize}
}
}

\subsection{Template para Segmentação}

\spec{{
\textbf{Template: Segmentação Odontológica} \\
\textbf{Problema:} Segmentar [estrutura anatômica] em [tipo de imagem]. \\
\textbf{Entrada:} \texttt{ndarray (H, W) uint8} — [descrição da imagem]. \\
\textbf{Saída:} \texttt{ndarray (H, W) int32} — máscara de segmentação. \\
\textbf{Critérios:}
\begin{itemize}
  \item dice\_score $>$ [0.75--0.90]
  \item iou $>$ [0.65--0.85]
  \item hausdorff\_95 $<$ [10--25] mm
\end{itemize}
}
}

\subsection{Template para Detecção}

\spec{{
\textbf{Template: Detecção Odontológica} \\
\textbf{Problema:} Detectar [condição clínica] em [tipo de imagem]. \\
\textbf{Entrada:} \texttt{ndarray (H, W) uint8} — [descrição da imagem]. \\
\textbf{Saída:} \texttt{list[dict]} — lista de detecções, cada uma com
\texttt{\{"bbox": (x,y,w,h), "confidence": float, "class": str\}}. \\
\textbf{Critérios:}
\begin{itemize}
  \item mAP@0.5 $>$ [0.70--0.85]
  \item recall $>$ [0.80--0.90]
  \item precision $>$ [0.75--0.88]
\end{itemize}
}
}

\section{Exemplos Completos de Especificação}

\subsection{Exemplo 1: Detector de Cáries (Classificação)}

\spec{
  \textbf{Sistema:} \texttt{detect\_caries\_bitewing} \\
  \textbf{Problema:} Detectar lesões cariosas proximais em radiografias
  bitewing digitais de dentes posteriores permanentes. A detecção segue
  o sistema ICDAS (International Caries Detection and Assessment System),
  classificando em escores de 0 (hígido) a 6 (cavidade extensa). \\
  \textbf{Entrada:} \texttt{ndarray (224, 224) uint8} — radiografia
  bitewing recortada na região interproximal, com equalização CLAHE
  aplicada e normalizada para intervalo [0, 255]. \\
  \textbf{Saída:} \texttt{dict} contendo:
  \begin{itemize}
    \item \texttt{"icdas\_score"}: int $\in$ [0, 6] — escore ICDAS predito
    \item \texttt{"has\_caries"}: bool — True se ICDAS $\geq$ 2
    \item \texttt{"confidence"}: float $\in$ [0, 1] — confiança da predição
    \item \texttt{"probabilities"}: ndarray (7,) — probabilidades por classe
  \end{itemize}
  \textbf{Critérios de aceitação:}
  \begin{itemize}
    \item accuracy $>$ 0,88
    \item sensitivity $>$ 0,85 (para ICDAS $\geq$ 2)
    \item specificity $>$ 0,85
    \item $\kappa$ de Cohen $>$ 0,80
    \item AUC-ROC $>$ 0,92
    \item tempo de inferência $<$ 150 ms (em CPU)
  \end{itemize}
}

\teste{Especificação aprovada pelo SpecVerifier. Campos obrigatórios
preenchidos, critérios numéricos verificáveis, tipos documentados.}

\subsection{Exemplo 2: Classificador de Lesões Bucais}

\spec{
  \textbf{Sistema:} \texttt{classify\_oral\_lesion} \\
  \textbf{Problema:} Classificar imagens intraorais de lesões de mucosa
  em 5 categorias clínicas: saudável, leucoplasia, líquen plano oral,
  carcinoma espinocelular (CEC) e úlcera aftosa. \\
  \textbf{Entrada:} \texttt{ndarray (224, 224, 3) uint8} — foto clínica
  intraoral RGB com iluminação padronizada e fundo removido. \\
  \textbf{Saída:} \texttt{dict} contendo:
  \begin{itemize}
    \item \texttt{"class"}: str — nome da classe predita
    \item \texttt{"probabilities"}: dict — \{classe: probabilidade\}
    \item \texttt{"heatmap"}: ndarray (224, 224) — mapa Grad-CAM
  \end{itemize}
  \textbf{Critérios de aceitação:}
  \begin{itemize}
    \item accuracy $>$ 0,90
    \item sensitivity $>$ 0,87 (para classes patológicas)
    \item specificity $>$ 0,87
    \item soma das probabilidades = 1,0 $\pm$ 1e-6
    \item AUC $\geq$ 0,93 para classe CEC (a mais crítica)
  \end{itemize}
}

\subsection{Exemplo 3: Segmentador de Dentes}

\spec{
  \textbf{Sistema:} \texttt{segment\_teeth\_panoramic} \\
  \textbf{Problema:} Segmentar dentes individuais em radiografias
  panorâmicas digitais, atribuindo a cada pixel um identificador de
  dente (1--32, conforme notação FDI) ou 0 para \textit{background}. \\
  \textbf{Entrada:} \texttt{ndarray (512, 512) uint8} — radiografia
  panorâmica em escala de cinza, redimensionada mantendo proporção. \\
  \textbf{Saída:} \texttt{ndarray (512, 512) int32} — máscara de
  segmentação com valores 0 (fundo) a 32 (dentes). \\
  \textbf{Critérios de aceitação:}
  \begin{itemize}
    \item Dice score $>$ 0,85 (média sobre 32 classes)
    \item IoU (Jaccard) $>$ 0,78
    \item Precisão $>$ 0,88
    \item Revocação $>$ 0,82
    \item Distância de Hausdorff (95º percentil) $<$ 12 mm
  \end{itemize}
}

\teste{Verificação do SpecVerifier: 3/3 especificações aprovadas.
Critérios formalmente parseados e verificáveis.}

\section{Usando o Decorator \texttt{@spec} do Pacote OdontoIA}

O coração do SDD no ecossistema OdontoIA é o decorador \texttt{@spec},
implementado no módulo \texttt{odontoia.sdd.spec}. Este decorador anexa
metadados de especificação a qualquer função Python, registrando-os em
um registro global consultável.

\begin{pratica}{Decorando funções odontológicas com @spec}
  Objetivo: Aplicar o decorador \texttt{@spec} a três funções
  representativas de problemas odontológicos reais e inspecionar
  as especificações registradas.
\end{pratica}

\codigo{codigos/cap12/spec_caries_detector.py}

\subsection{Anatomia do @spec}

O decorador aceita cinco parâmetros:

\begin{description}
  \item[\texttt{problem}] (str, obrigatório): descrição do problema
    odontológico em linguagem natural.
  \item[\texttt{input}] (str, obrigatório): especificação do tipo e
    forma da entrada, preferencialmente com tipos NumPy/PyTorch.
  \item[\texttt{output}] (str, obrigatório): especificação do tipo e
    semântica da saída.
  \item[\texttt{criteria}] (List[str], obrigatório): lista de critérios
    de aceitação, preferencialmente com operadores relacionais
    (\texttt{>}, \texttt{<}, \texttt{>=}, \texttt{<=}, \texttt{==}).
  \item[\texttt{category}] (str, opcional): categoria para agrupamento
    (\texttt{"classificacao"}, \texttt{"segmentacao"},
    \texttt{"deteccao"}, \texttt{"preprocessamento"}).
\end{description}

\subsection{Recuperando e Validando Especificações}

Uma vez registrada, a especificação pode ser recuperada e validada:

\begin{itemize}
  \item \texttt{get\_spec(func)} — retorna o dicionário de especificação
    associado à função, ou \texttt{None} se não especificada.
  \item \texttt{validate\_spec(func)} — valida se todos os campos
    obrigatórios estão presentes e se os critérios são bem formados.
  \item \texttt{list\_specs(category=None)} — lista todas as funções
    registradas, opcionalmente filtradas por categoria.
\end{itemize}

\subsection{SpecVerifier: Verificação Automatizada}

O módulo \texttt{odontoia.sdd.validator} fornece a classe
\texttt{SpecVerifier}, que executa verificação programática:

\begin{itemize}
  \item \texttt{SpecVerifier(strict\_mode=False)} — instancia o
    verificador; \texttt{strict\_mode=True} levanta exceção na
    primeira falha.
  \item \texttt{verify(func, *args, test\_output=None)} — executa a
    função com os argumentos fornecidos e valida contra a especificação.
  \item \texttt{summary()} — retorna um resumo de todas as verificações
    realizadas (total, aprovadas, reprovadas, taxa de sucesso).
\end{itemize}

\teste{Ao executar \texttt{SpecVerifier().verify(detect\_caries\_bitewing)},
a especificação é validada com sucesso: campos obrigatórios presentes,
critérios parseáveis, função documentada.}

\testefalha{Ao executar \texttt{SpecVerifier(strict\_mode=True).verify(func\_sem\_spec)},
o SpecVerifier levanta \texttt{SpecError}: ``Função não possui especificação @spec''.}

\section{Síntese do Capítulo}

Neste capítulo, estabelecemos os fundamentos do Desenvolvimento Orientado
por Especificação (SDD) aplicado à Odontologia. Os principais pontos
a reter são:

\begin{itemize}
  \item \textbf{SDD é a disciplina de especificar antes de implementar},
    garantindo que cada função tenha um contrato formal verificável.

  \item \textbf{A Odontologia é um domínio \textit{safety-critical}}:
    erros de software podem causar danos reais a pacientes, tornando o
    SDD uma exigência ética, não apenas técnica.

  \item \textbf{Uma boa especificação tem 4 componentes}: problema,
    entrada, saída e critérios de aceitação verificáveis — todos
    explícitos e tipados.

  \item \textbf{O ciclo SDD tem 4 fases}: Spec $\rightarrow$ Implement
    $\rightarrow$ Verify $\rightarrow$ Deploy, com refinamento contínuo.

  \item \textbf{O pacote \texttt{odontoia.sdd}} fornece o decorador
    \texttt{@spec} e o verificador \texttt{SpecVerifier}, automatizando
    todo o processo.

  \item \textbf{Templates de especificação} aceleram o trabalho para
    classificação, segmentação e detecção — as três tarefas mais comuns
    em IA odontológica.
\end{itemize}

No próximo capítulo, integraremos o SDD com o TDD (\textit{Test-Driven
Development}), completando o ciclo de qualidade: especificar o que fazer
(SDD) e testar se foi feito corretamente (TDD).
